\section{The Positional Price of Keeping the Tree}\label{sec:bounds}
The positional cost comes from two sources: constraints imposed by the hierarchy and additional choices made by a layout. We first isolate the hierarchy in a point model, then account for interval packing, and finally allow controlled changes to merge levels. This progression answers RQ1--RQ3 using the same reference coordinate. Suppress $t$, write $[v]^+=\max(v,0)$ and take empty maxima as zero; complete packing derivations are in the supplementary material.

\subsection{Exact Topology}
The first question removes feature width and asks what the hierarchy alone forces. Exact topology keeps sibling subtrees in separate blocks, so their reference hulls determine the cost of choosing either block orientation. Let $[m_X,M_X]$ be the reference hull of subtree $X$ in the reduced equal-level tree.
\begin{theorem}[Exact-topology position bound]\label{thm:order}
Every exact profile satisfies $\nroe(\pi)\ge\nu^*$, where
\begin{equation}\label{eq:nu}
\nu^*=\frac{1}{2L}\max_v\max_{X\ne Y\in\mathrm{ch}(v)}[\min(M_X-m_Y,M_Y-m_X)]^+.
\end{equation}
Recursive ordering by child-hull midpoint attains $\nu^*$ in $O(n\log n)$ time.
\end{theorem}
Sibling blocks must take one orientation. Its inversion cost is either $[M_X-m_Y]^+/2$ or $[M_Y-m_X]^+/2$, and midpoint order chooses the smaller. Every inverted leaf pair separates at its LCA. Therefore error above $\nu^*$ can be removed by reordering in the point model.

\subsection{Finite Widths}
Set $H_i=\rho W_i/2$, $\ell_i=q_i-W_i/2-H_i$, $h_i=q_i+W_i/2+H_i$ and $s_i=W_i+g$. Packing summaries
\begin{equation}\label{eq:packing}
\begin{split}
A(S)&=\max_{a\in S}\bigl(\ell_a+\sum_{i\in S:\ell_i\ge\ell_a}s_i\bigr),\\
B(S)&=\max_{b\in S}\bigl(\sum_{i\in S:h_i\le h_b}s_i-h_b\bigr)
\end{split}
\end{equation}
measure the interval span forced by upper and lower reference-endpoint sets. Define $\omega(X,Y)=\tfrac12\min\{A(X)+B(Y),A(Y)+B(X)\}-g/2$, and $E_\partial=\max\{0,B(\mathcal F)-g,A(\mathcal F)-g-L\}$ for all leaves $\mathcal F$.
\begin{theorem}[Width-aware position bound]\label{thm:width}
For every feasible exact-topology interval layout,
\begin{equation}\label{eq:widthbound}
\|x-q\|_\infty\ge\tau(\pi)\ge\mathcal B_w
=\max\{E_\partial,\max_{(X,Y)\in\mathcal P_0}\omega(X,Y)\},
\end{equation}
where $\mathcal P_0$ contains distinct sibling descendant-leaf sets. With zero widths and gaps, $\mathcal B_w=L\nu^*$.
\end{theorem}
The first and last anchors must accommodate their intervals and intervening gaps up to the permitted anchor slack. Sibling contiguity chooses a block orientation, while the canvas ends supply $E_\partial$. Thus ordering alone no longer determines the positional cost: even an order attaining the point bound may require additional displacement to fit its intervals. The bound has coordinate units, while $\nu^*$ is normalized by $L$. This motivates evaluating the exact packing cost of each candidate order before optimizing its anchors.

\subsection{Bounded Merge Error}
If leaf $k$ lies between descendants $i,j$ of node $u$, interval containment gives $V_{ik}^{\psi}\le V_{ij}^{\psi}$, hence $V_{ik}\le h(u)+2\eta$. Let $A_\eta(u)$ be the highest ancestor with $h(A_\eta(u))-h(u)\le2\eta$. Form $\mathcal P_\eta$ from pairs $(\mathcal D(u),\{k\})$ with $k\notin\mathcal D(A_\eta(u))$, and pairs $(\mathcal D(u),\mathcal D(w))$ with each set disjoint from the other's $A_\eta$ ancestor set, for internal $u,w$.
\begin{theorem}[Bounded-topology position bound]\label{thm:chain}
If $\dtop\le\eta$, every pair in $\mathcal P_\eta$ lies wholly in one block orientation. Thus
\begin{equation}\label{eq:chainbounds}
\begin{split}
\nroe(\pi)&\ge C_\eta=\max_{(X,Y)\in\mathcal P_\eta}\frac{[\min(M_X-m_Y,M_Y-m_X)]^+}{2L},\\
\tau(\pi)&\ge C_{w,\eta}=\max\{E_\partial,\max_{(X,Y)\in\mathcal P_\eta}\omega(X,Y)\}.
\end{split}
\end{equation}
\end{theorem}
The accumulated ancestor gap determines which outside leaves may enter a block, so an allowance releases only the conflicts supported by sufficiently small merge-level differences. Flattening parent gaps of at most $2\eta$ gives a second certificate; we use the larger of the two valid bounds. Together, the three results distinguish forced displacement from displacement chosen by a layout. Their algorithmic implication is to search orders using the exact fixed-order cost $\tau$, admit them only when their merge error fits the allowance, and limit additional displacement during anchor placement. A certificate is a lower bound; the construction must still produce and check a feasible map.
